\chapter{Verzeichnis der Hilfsmittel}\label{sec:AnhangHilfsmittel}
\section{Programme}


\begin{tabular}{|p{4cm}|p{12cm}|}
	\hline 
Textverarbeitung	&  PDF Latex , TeX Live (\url{https://tug.org/texlive}) version 2022\\ 
	\hline 
Grafik	&  Corel Draw V12\\ 
	\hline 
Literaturverwaltung	& JabRef 5.11 \\ 
	\hline 
	&  \\ 
	\hline 
	&  \\ 
	\hline 
	&  \\ 
	\hline 
\end{tabular} 



\section{KI Werkzeuge}\label{sec:AnhangHilfsmittel_KI}
\begin{tabular}{|p{4cm}|p{12cm}|}
	\hline 
Titelbild	& erzeugt unter Verwendung von  \url{https://www.bing.com/images/create} \\ 
\hline 
Bild \ref{fig:beispiel1} rechts & erzeugt unter Verwendung von  \url{https://www.bing.com/images/create} \\ 
	\hline 
Kapitel \ref{sect:KItools}	& \url{https://chat.mistral.ai/}: Text mit Vorschl"agen generiert und einige Textpassagen umformuliert  \\ 
	\hline 
	&  \\ 
	\hline 
	&  \\ 
	\hline 
\end{tabular} 



\chapter{Latex Nutzung}
\section{Latex DO's \& DON'Ts}
						
Zum richtigen Umgang mit \LaTeX wird die PDF-Datei "`l2tabu"' empfohlen. In ihr sind die wichtigsten Regeln im Umgang mit \LaTeX und dem Koma-Script zusammengefasst (\url{www.mast.queensu.ca/~andrew/LaTeX/latex-dos-and-donts.pdf}).




\section{Grafiken mit TikZ und PGFPlots}

Eine typische Aufgabe in Abschlussarbeiten ist die Erstellung von Grafiken. Eine Möglichkeit ist die Verwendung eines speziellen Zeichenprogramms, wie zum Beispiel Inkscape oder CorelDraw. Aber gerade bei einfachen Blockschaltbildern können andere Lösungen auch überzeugen.

Ein Beispiel ist das Programmpaket TikZ. Hier können in einer Beschreibungssprache solche Grafiken komponiert werden.
Die Verwendung ist nicht ganz einfach, aber Ruben Sander hat mit seinem Betreuer Alfred Hülkenberg eine schöne Seminararbeit verfasst, die die Verwendung gut erklärt \cite{Sander2021}.

Eine weitere praktische Erweiterung dazu ist PgfPlot, die TikZ um die Möglichkeit erweitert schöne Plots aus Messdaten zu erzeugen. Es gibt sogar Tools, dies direkt aus Matlab oder Python zu realisieren. Hier hat Julian Bigalke mit seinem Betreuer Alfred Hülkenberg gleichermaßen eine sehr informative Seminararbeit verfasst \cite{Bigalke2023}.

Beide Manuale werden in diesem Vorlagenpaket im Zip Archiv mitgeliefert. 

\cite{abcde}